\section{Background}

% \subsection{Dec-POMDP}
A fully cooperative multi-agent task can be formulated as a Dec-POMDP~\cite{oliehoek2016concise}, which is defined as a tuple $\mathcal{G}=\langle N, S, A, P, R, O, \Omega, n, \gamma\rangle$, where $N$ is a finite set of $n$ agents, $s\in S$ is the true state of the environment, $A$ is the set of actions, and $\gamma\in [0, 1)$ is a discount factor. At each time step, each agent $i\in N$ receives his own observation $o_i\in \Omega$ according to the observation function $O(s, i)$, and selects an action $a_i\in A$, which results in a joint action vector $\bm{a}$. The environment then transitions to a new state $s'$ based on the transition function $P(s'| s,\bm{a})$, and inducing a global reward $r=R(s, \bm{a})$ shared by all the agents. Each agent has its own action-observation history $\tau_i\in \mathcal{T}_i \doteq (\Omega_i\times A)^*$. Due to partial observability, each agent conditions its policy $\pi_i(a_i | \tau_i)$ on $\tau_i$. The joint policy $\bm{\pi}$ induces the joint action-value function $Q_{tot}^{\bm{\pi}}(s, \bm{a})=\mathbb{E}_{s_{0:\infty},\bm{a}_{0:\infty}}\left[\sum_{t=0}^\infty \gamma^t r_t \mid s_0=s, \bm{a}_0=\bm{a}, \bm{\pi} \right]$.
% the joint value function $V^{\bm{\pi}}(s) = \mathbb{E}_{s_{0:\infty},\bm{a}_{0:\infty}}\left[\sum_{t=0}^\infty \gamma^t r_t| s_0=s, \bm{\pi} \right]$, and 

\subsection{Centralized Training with Decentralized Execution}
Our method adopts the framework of centralized training with decentralized execution (CTDE)~\cite{lowe2017multi, foerster2018counterfactual, sunehag2018value, rashid2018qmix, son2019qtran, son2020qtran, wang2020qplex, ma2021modeling}. This framework tackles the exponentially growing joint action space by decentralizing the control policies while adopting centralized training to learn cooperation. Agents learn in a centralized manner with access to global information but execute based on their local action-observation history. One promising approach to implement the CTDE framework is value function factorization. The IGM (individual-global-max) principle~\cite{son2019qtran} guarantees the consistency between the local and global greedy actions. When IGM is satisfied, agents can obtain the optimal global action by simply choosing the local greedy action that maximizes each agent's individual utility function $Q_i$. Some algorithms have successfully used the IGM principle~\cite{rashid2018qmix,wang2020qplex,rashid2020weighted} to push forward the progress of MARL.

%This principle are wildly used ~\cite{rashid2018qmix,wang2020qplex,rashid2020weighted} to push forward the progress of MARL.

%Some algorithms have successfully used the IGM principle~\cite{rashid2018qmix,wang2020qplex,rashid2020weighted} to push forward the progress of MARL.
% \begin{equation}
% 	\argmax_{\bm{a}\in A^n} Q_{tot}(\bm{\tau}, \bm{a}) = \left(
% 	\begin{aligned}
% 		&\argmax_{a_1\in A}Q_1(\tau_1, a_1)\\
% 		&\qquad \quad\ldots\\
% 		&\argmax_{a_n\in A}Q_n(\tau_n, a_n)\\
% 	\end{aligned}
% 	 \right),
% \end{equation}
